%% Does AI rate through the cycle? — Finance Research Letters (Elsevier CAS single-column)
\documentclass[a4paper,fleqn]{cas-sc}
\usepackage[T1]{fontenc} % scalable Type-1 fonts (fixes Type-3/ligature extraction)
\usepackage{lmodern}
\usepackage{microtype}
\usepackage[authoryear,longnamesfirst]{natbib}
\usepackage{graphicx}
\usepackage{booktabs}
\usepackage{verbatim}
\usepackage{url}

% listings with breaklines: the verbatim prompts in Appendix B are single logical
% lines of up to 615 characters, which verbatim ran off the page and clipped.
\usepackage{listings}
\lstset{breaklines=true,breakatwhitespace=false,basicstyle=\ttfamily\scriptsize,
 columns=fullflexible,frame=none,xleftmargin=0pt,upquote=true,
 showstringspaces=false,breakindent=0pt,postbreak=}
\begin{document}
\let\WriteBookmarks\relax
\def\floatpagepagefraction{1}
\def\textpagefraction{.001}

\shorttitle{Does AI rate through the cycle?}
\shortauthors{S. Chincholikar and R. Chawla}

\title[mode=title]{Does AI rate through the cycle? Procyclicality in credit assessment}

\author[1]{Samir Chincholikar}[orcid=0009-0007-2779-3492]
\ead{samir.chincholikar@gmail.com}
\credit{Conceptualization, Methodology, Software, Formal analysis, Writing -- original draft}
\affiliation[1]{organization={Independent researcher}, city={New York}, country={USA}}

\author[2]{Robin Chawla}[orcid=0009-0007-2807-3948]
\cormark[1]
\ead{robin.chawla.cse14@iitbhu.ac.in}
\credit{Conceptualization, Validation, Writing -- review and editing, Supervision}
\affiliation[2]{organization={Independent researcher}, city={New York}, country={USA}}

\cortext[1]{Corresponding author}

\begin{abstract}
Large language models (LLMs) are increasingly proposed as credit-rating and default-assessment
tools, yet whether they rate ``through the cycle'' or amplify it has been asserted rather than
measured. We hold
synthetic firm fundamentals byte-identical across an ordered five-point macroeconomic-severity axis
(plus a matched placebo axis) and elicit letter ratings and one-year default
probabilities from seven models spanning two families from the lowest cost tier to the capability
frontier, so that any systematic rating movement along the macro axis is procyclicality by construction, identified without the
unobserved-fundamentals confound that limits the classical Amato--Furfine test. Every model is
significantly procyclical, from $0.182$ to $0.775$ rating notches of downgrade per severity step, and
$18$ of $21$ model pairs differ significantly. Capability lowers the level but not the character: neither
frontier model differs significantly from a lower cost Flash-tier model, and downside
sensitivity exceeds upside in every model. An explicit ``through-the-cycle'' instruction
does not restore through-the-cycle behaviour. It anchors the rating while suppressing the
point-in-time default probability that should keep moving, and it fails hardest at the frontier, where
the rating flattens to zero or below and so does default-probability retention. A vaguer
``be stable'' instruction preserves the distinction far better, pointing to cycle-specific wording, not
stabilisation as such. The pattern survives on real SEC-filer fundamentals, and the
capital consequence keeps its sign across every mapping we vary. Among the tested models,
greater capability does not eliminate this channel: it produces the most stable-looking rating while carrying
the least risk signal.
\end{abstract}

\begin{highlights}
\item First controlled test of procyclicality in LLM credit ratings.
\item All seven models downgrade identical firms by 0.18 to 0.78 notches per severity step.
\item Frontier models do not differ significantly from a lower cost Flash model.
\item The effect is asymmetric: downside sensitivity exceeds upside in every model.
\item A ``through-the-cycle'' prompt anchors ratings rather than de-cyclifying them.
\end{highlights}

\begin{keywords}
large language models \sep credit ratings \sep procyclicality \sep through-the-cycle \sep
prompt framing \sep financial regulation
\end{keywords}

\maketitle

\section{Introduction}\label{sec:intro}
Credit ratings are important for debt pricing, investment decisions, and regulatory capital
requirements. A perennial concern is that ratings could be procyclical, with companies downgraded
during recessions and upgraded during expansions even when the underlying credit fundamentals have not
materially changed. Such rating changes can exacerbate the credit cycle by raising borrowing costs and,
through internal-ratings-based (IRB) capital rules, capital requirements when economic conditions are
already weak \citep{kashyap2004,gordy2006}. Rating agencies therefore attempt to capture
creditworthiness ``through the cycle'' (TTC), placing more importance on a borrower's ability to
withstand temporary macroeconomic fluctuations than on mechanically reacting to current conditions
\citep{amato2004,loffler2004}. However, whether these ratings actually serve this purpose remains
debated \citep{loffler2013}.

There is growing interest in using large language models (LLMs) for financial decision-making,
including credit decisioning and risk assessment, and supervisors have flagged that generative models
could re-introduce procyclicality \citep{bis2024}. Institutional use is already emerging: Moody's offers AI-powered
systems for company credit assessment, underwriting, and credit-memo generation, with documented
adoption by banks such as Taishin International Bank, while S\&P Global has launched an agentic-AI
Credit Memo Builder for loan committees, underwriters, and credit analysts
\citep{moodys2026a,moodys2026b,spglobal2026}. These systems retain analyst oversight rather than
replacing final credit judgment, but their deployment makes LLM behaviour directly relevant to
these workflows. LLMs behave as
context-sensitive economic agents \citep{horton2023} and can exhibit framing sensitivity
\citep{binz2023}, echoing classical framing effects \citep{tversky1981}, and a shared foundation
model could homogenise decisions across users \citep{bommasani2022}. Yet the size of this effect on credit ratings has not
been directly measured. In observational studies of human agencies, fundamentals and macroeconomic
conditions move together, so the classical test of \citet{amato2004} must proxy for unobserved
fundamentals.

This paper uses a controlled experimental setting to test the procyclicality of LLM-based credit
rating. We keep the firm fundamentals unchanged and only change the macroeconomic environment, so any
systematic change in the assigned rating can be attributed to the macroeconomic context and not to
changes in firm credit quality. The benchmark is not that any movement is bad: a genuine
through-the-cycle rater should hold the \emph{rating} steady while its point-in-time default
probability still tracks conditions, and we test that joint condition directly. We make three
contributions. First, all seven models, including the frontier tier of both families, downgrade
identical firms as conditions deteriorate. The level falls with capability, but neither frontier model
differs significantly from a Flash-tier model. Second, the response is strongly
asymmetric, with downside sensitivity exceeding upside sensitivity in every model. Third, prompt design
does not fix it: an explicit through-the-cycle instruction flattens the rating while suppressing the
one-year default probability's responsiveness, so the model anchors both outputs rather than separating
point-in-time risk from a through-the-cycle rating, echoing classical anchoring effects
\citep{tversky1974}, and it does this most completely at the frontier. 

\section{Related literature}\label{sec:lit}
Our study relates to three strands. The first assesses through-the-cycle behaviour and rating
procyclicality: whether ratings respond to the business cycle once fundamentals are controlled for,
and whether agencies in fact achieve through-the-cycle stability
\citep{amato2004,loffler2004,loffler2013}. Rating levels also move market pricing directly
\citep{cantor1996}. The prudential stakes come from the interaction of ratings
with Basel capital rules \citep{kashyap2004,gordy2006}.

The second strand is the rapid adoption of \emph{LLMs in finance}
\citep{ali2025}, spanning return prediction, financial-statement analysis, financial advice and
simulated economic agents \citep{lopezlira2023,kim2024,fieberg2025,horton2023}
and, most relevant here, reliability evidence for banking supervision \citep{garcia2026},
with related work measuring policy uncertainty and AI-firm financing
\citep{ito2025,giaretta2026}. Most closely, \citet{hens2025} compare the risk profiles LLMs produce for investors
with those of experienced bankers, finding relatively small differences in the levels of risk
assigned, although the models give weaker explanations. We examine a different dimension of reliability: whether an LLM alters its credit assessment in response to rating-irrelevant
changes in the macroeconomic context surrounding otherwise identical firm information.

Anchoring, framing and elicitation-dependent calibration are documented in language-model decisions
\citep{binz2023,kadavath2022}, echoing the classical human framing literature
\citep{tversky1974,tversky1981}, and shared foundation models can homogenise decisions across
users \citep{bommasani2022}. In our setting these behavioural properties
connect to a specific financial mechanism: if the macroeconomic context affects credit ratings while firm
fundamentals are held fixed, contextual sensitivity translates directly into rating procyclicality.

\section{Design and identification}\label{sec:design}
The unit of analysis is a rating cell $=$ (model, prompt framing, seed, firm, macro state). We
construct $80$ synthetic firms spanning five credit-quality tiers and eight sectors, and characterise
each firm using the same set of credit fundamentals: revenue growth, EBITDA margin, leverage, interest
coverage, free cash flow, liquidity, and Altman-style financial ratios \citep{altman1968}. The
information about each firm presented to the models is the same in all experimental conditions. Each
firm's fundamentals are serialised once and reused byte-identically under every macro state. Following
this, we merely alter the model's exposure to the macroeconomic environment by switching it between
five ordered states: boom, expansion, neutral, slowdown, and severe recession (severity
$-2\ldots+2$). A common set of macroeconomic indicators is used to represent each state, including GDP
growth, unemployment, credit spreads, default rates, and lending standards. No calendar dates are
given, reducing the chance that a model associates a situation with a particular historical event.

To strengthen identification, we add two controls. We first generate a matched placebo condition
carrying credit-irrelevant information about regional weather and traffic. The placebo matches the macro scenarios in format and numerical
density but is credit-irrelevant, so it separates macroeconomic sensitivity from a general tendency to
respond to changing context. The net (macro minus placebo) coefficient is the demand-effect-adjusted
cyclicality. Second, a firm and framing keep the same seed across all macro states, so sampling
variation is not systematically related to the macro state. Each assessment produces a letter credit rating (mapped to an
S\&P notch on a 21-point scale) and a one-year probability of default (PD).

We consider four prompt framings: a terse baseline, a point-in-time instruction, an explicit
through-the-cycle instruction, and a generic ``be stable'' instruction that makes no reference to the
credit cycle. The experiment uses seven models spanning the full capability range of two families: four from OpenAI
(\texttt{gpt-4o-mini}, \texttt{gpt-4.1-mini}, \texttt{gpt-4o}, \texttt{gpt-6-astra}) and three Google
(\texttt{gemini-flash-lite-latest}, \texttt{gemini-flash-latest}, \texttt{gemini-pro-latest}). Three OpenAI arms call dated snapshots. The other four use undated aliases that cannot be version-resolved (Appendix~\ref{app:robust}). Each cell is run at temperature zero where supported, which minimises decoding variability and isolates the macro axis as the treatment contrast, with replicate seeds: the primary condition carries five draws per cell for the four
models whose additional-seed captures completed (\texttt{gpt-4o-mini}, \texttt{gpt-4.1-mini},
\texttt{gpt-4o}, \texttt{gemini-flash-latest}) and two for the remaining three
(\texttt{gemini-flash-lite-latest}, \texttt{gemini-pro-latest}, \texttt{gpt-6-astra}), whose
extra-seed captures exceeded the pre-registered $2$\% error-rate gate, which refuses a whole
capture rather than individual observations. Two frontier-model properties constrain what can be held fixed: \texttt{gpt-6-astra} refuses temperature zero and runs at its default, and both frontier
models reason internally by construction. Because decoding randomness is unrelated to macro
severity, it should widen intervals rather than generate a systematic slope. \texttt{gemini-pro} does accept temperature zero, so the frontier result is reproduced at the
design's own decoding setting.
The dependent variable is the change in rating notch as severity rises,
positive meaning more procyclical. Intervals are $2{,}000$-draw firm-clustered, tier-stratified bootstraps. Tables~\ref{tab:permodel} and~\ref{tab:robust} and Fig.~\ref{fig:asym} report the secondary estimands: swing,
half-slopes, the Amato--Furfine ordered probit and the log-probability expected-notch slope
where the provider exposes it. 

\section{Results}\label{sec:results}
\subsection{LLMs are procyclical on identical fundamentals}\label{sec:pro}
On byte-identical fundamentals, all seven models downgrade the same firms as the described macroeconomy
worsens, and no model's interval contains zero (Table~\ref{tab:permodel} and Fig.~\ref{fig:main}a). At the pooled terse
coefficient (Table~\ref{tab:gradient}) a BBB firm moves toward BBB$-$/BB$+$ with nothing real changing.

Model choice is a first-order determinant: under a paired firm-clustered bootstrap $18$ of the $21$
model pairs differ significantly, so the fourfold spread in Table~\ref{tab:permodel} is tested rather
than inferred from overlapping intervals. Its structure matters more than its width. Within each family
the coefficient declines monotonically in point estimate with capability, yet capability does not order it \emph{across}
families: the two tested frontier models converge near $0.19$ from very different starting points
rather than approaching zero. Before correction the only three pairs that do \emph{not} differ are Gemini Flash versus Gemini Pro,
Flash versus GPT-6 Astra, and Pro versus Astra: neither frontier model differs
significantly from a Flash-tier model. Capability buys a lower level, not a different character,
and Section~\ref{sec:ttc} shows the character worsens.

The result survives the checks in Table~\ref{tab:robust}: the credit-irrelevant placebo moves
ratings only $0.0425$ notches per step across the five original models, leaving a net effect of $0.387$ (95\% CI $[0.363,0.412]$) that
still excludes zero, and the coefficient holds under served-model fixed effects, a log-probability
estimator and the Amato--Furfine ordered-probit estimator.

Magnitude is therefore a model-selection variable: the swings in Table~\ref{tab:permodel} span over
fourfold, frontier models under one notch and GPT-4o-mini above three. The claims we press are the
structural ones, which do \emph{not} attenuate with capability ($\S$\ref{sec:ttc}).

\subsection{The procyclicality is asymmetric}\label{sec:asym}
Split at neutral, downgrade sensitivity is about $3.2$ times upgrade sensitivity pooled, and every
model's downside-to-upside interval excludes~$1$ (Table~\ref{tab:permodel} and Fig.~\ref{fig:asym}). That is a direction, not an
ordering: the ratio's denominator is small, so the per-model ranking is not resolvable. The direction is
what matters, because downgrades in downturns raise borrowing costs and capital requirements exactly
when conditions are already strained. Because capital requirements are convex in the rating and bind hardest in stress,
an asymmetric downgrade profile is more procyclical than a symmetric one of equal average slope.

\subsection{``Through-the-cycle'' instructions anchor rather than de-cyclify}\label{sec:ttc}
Prompt framing is a strong determinant of rating sensitivity: the pooled slope falls monotonically
across the four framings, roughly threefold from point-in-time to explicit through-the-cycle
(Table~\ref{tab:gradient}). At first sight that suggests through-the-cycle prompting reduces
procyclicality. It does not. Genuine through-the-cycle behaviour is a \emph{joint} condition: the rating slope should
fall toward zero \emph{and} the point-in-time default probability should keep moving with the cycle.
Neither is required to be exact, so we operationalise both as tolerance bands tested against
their intervals: a rating slope within $\pm0.05$ notches per step of zero, and PD retention of at
least $0.75$ (Appendix~\ref{app:robust}). We measure the default probability on the
log-odds scale, which reduces dependence on PD \emph{level} differences between models. That is not an external calibration test, and the calibration of verbalised LLM probabilities is task- and elicitation-dependent \citep{kadavath2022}. We therefore do not rest the anchoring result on the
elicited PDs alone. It is supported jointly by three measures: the rating-token log-probability slope,
which never reads a stated PD and falls comparably under the instruction wherever the provider exposes it (Table~\ref{tab:robust}), PD retention, which uses the elicited one-year PDs on
the log-odds scale, and the ``be stable'' versus through-the-cycle control, which holds the elicitation fixed and varies
only the instruction.

No model satisfies both conditions, and they fail in opposite directions
(Table~\ref{tab:permodel}). The five smaller models fail the rating condition while retaining part of
the PD slope. The frontier models invert that: Gemini Pro alone has a residual interval inside the band and retains
none of the PD slope, while GPT-6 Astra overshoots it, not merely ceasing to track the cycle but
reversing. Neither is through-the-cycle behaviour. The instruction stabilises the rating and suppresses the
default probability that should remain sensitive: it anchors the wrong object, and most completely
where capability is highest. The vaguer ``be stable'' instruction, which is not semantically about the
cycle, leaves the rating slope higher while preserving far more default-probability
movement, consistent with the additional suppression arising from cycle-specific wording rather than from stabilisation as such.

\subsection{Capital consequences}\label{sec:capital}
What matters is not that the rating moves but what the movement costs. We translate
the rating path into capital by anchoring through the \emph{rating}, which is what enters a credit pipeline rather than the disclaimed verbalised PD: we map each notch to a stylised one-year default
probability loosely calibrated to published one-year corporate default rates
\citep{spratings2025}, and apply the Basel
internal-ratings-based (IRB) corporate risk-weight formula \citep{bcbs_cre31} as a proxy for the
economic logic (parameters in Appendix~\ref{app:robust}). Reporting the \emph{change} across the severity axis,
an identical firm's required capital rises, across the five original models, by a pooled $1.81$ percentage points of exposure from
boom to severe recession (95\% CI $[1.67,1.94]$), and by nearly twice that for the most sensitive model.
Because the risk weight is convex in the rating, the $\S$\ref{sec:asym} asymmetry compounds: the downside
increase is about five times the upside decrease, a wider gap than the notch asymmetry alone. A
$259$-cell sensitivity grid leaves the sign positive in every cell (Appendix~\ref{app:robust}), so the
direction is robust while the magnitude is assumption-dependent.

\section{Conclusion}\label{sec:concl}
With fundamentals fixed by construction, all seven models, including the frontier tier of both families, downgrade identical firms as conditions are described as worsening, and do so
asymmetrically, punishing simulated downturns far more than they reward simulated expansions. The
level of the effect falls with capability. Its structure does not. Neither frontier model differs
significantly from a Flash-tier model, so the choice of model remains a prudential
decision, but a more capable model is not an escape from the behaviour.

Prompt design mitigates it only partially, and at the frontier not in the way that matters. An
explicit through-the-cycle instruction flattens the rating to zero or below for both frontier models, while
driving the one-year default probability's cyclical sensitivity to zero or below. The models stabilise
both outputs instead of separating a through-the-cycle rating from point-in-time risk, and they do so
most completely where capability is highest: the frontier models produce the most stable-looking rating
while carrying the least risk signal. Translated into capital requirements, the direction of the swing
survives every mapping, loss-given-default and maturity assumption we vary, though its magnitude does
not.

What the experiment demonstrates is a behavioural mechanism in the served endpoints at the collection
epoch, not a property of a model family in perpetuity and not an evaluation of deployed credit systems: we observe no rating committee, no lending decision, and no
analyst acting on these outputs. Varying the macroeconomic context around fixed firm information does
capture an exposure that arises in real pipelines, where models reach macroeconomic information through
market data, news, retrieval or analyst inputs. Replacing the synthetic fundamentals with real
de-identified filing-derived fundamentals leaves the effect in place, which narrows the realism gap on the firm side but
not on the others, and the frontier model's ability to recognise roughly a sixth of those tested,
despite name, ticker and date removal, is itself why the primary design is synthetic. Future work could benchmark against agency and internal-rating baselines, and test
whether fine-tuning or retrieval design can reduce procyclicality without suppressing genuine changes
in credit risk.

\begin{table}[t]\centering\small
\caption{Per-model results. \textbf{Panel A}: procyclicality on byte-identical fundamentals. $\beta$ = notches of downgrade per $+1$ macro-severity step under the terse framing, with firm-clustered bootstrap 95\% CI; \textbf{Swing} = measured boom$\to$severe-recession notches (not $4\beta$: the endpoint ratio runs $3.96$--$4.39$, so the gap carries curvature); \textbf{D:U} = downside over upside half-slope. Fig.~\ref{fig:asym} plots the two half-slopes with their intervals. The ratio has its own clustered intervals, with lower bounds from $1.23$ to $3.06$ and upper bounds reaching $8.50$. All exclude one, and they overlap across models, so they do not rank them. \textbf{Resid.} and \textbf{PD-ret.} = $\beta$ and the retained fraction of the log-odds PD slope under the through-the-cycle instruction; \textbf{$\beta$ real} = the same coefficient on the real-fundamentals battery and \textbf{$\Delta\beta$} its drift on re-measurement 2.5 months later (blank where the model was not in that arm). \textbf{Panel B}: one macro indicator moved across the same $-2\ldots+2$ range with the other four held neutral, so each column is comparable to $\beta$ above; $\Sigma/\beta$ sums the five and divides by the joint slope, below one meaning the indicators reinforce one another. $^{\dagger}$interval includes zero.}
\label{tab:permodel}
\begin{tabular}{lrrrrrrr}
\toprule
\multicolumn{8}{l}{\emph{Panel A: procyclicality, through-the-cycle response, and revision arms}}\\
\midrule
Model & $\beta$ (95\% CI) & Swing & D:U & Resid. (95\% CI) & PD-ret. & $\beta$ real & $\Delta\beta$ \\
\midrule
Gemini Pro & 0.182 [0.15, 0.22] & 0.72 & 2.3 & $-$0.021 [$-$0.04, 0.00] & 0.00 &  &  \\
GPT-6 Astra & 0.198 [0.18, 0.22] & 0.87 & 4.8 & $-$0.055 [$-$0.08, $-$0.04] & $-$0.11 & 0.151 &  \\
Gemini Flash & 0.230 [0.19, 0.27] & 0.96 & 4.0 & 0.059 [0.03, 0.08] & 0.22 & 0.094 & $-$0.037 \\
Gemini Flash-Lite & 0.285 [0.25, 0.32] & 1.15 & 3.6 & 0.056 [0.03, 0.08] & 0.34 &  &  \\
GPT-4o & 0.387 [0.34, 0.43] & 1.55 & 3.3 & 0.187 [0.15, 0.22] & 0.52 & 0.407 & +0.002 \\
GPT-4.1-mini & 0.471 [0.41, 0.53] & 1.94 & 3.9 & 0.096 [0.06, 0.13] & 0.24 &  & $-$0.021 \\
GPT-4o-mini & 0.775 [0.69, 0.86] & 3.18 & 2.6 & 0.403 [0.34, 0.47] & 0.59 & 0.814 & $-$0.028 \\
\midrule
\multicolumn{8}{l}{\emph{Panel B: macro-factor decomposition (one indicator at a time)}}\\
\midrule
Model & GDP & Unem. & Spread & Default & Standards & $\Sigma/\beta$ & \\
\midrule
GPT-4o-mini & 0.184 & 0.197 & 0.148 & 0.199 & 0.175 & 1.17 &  \\
GPT-4o & 0.019 & 0.043 & 0.112 & 0.114 & 0.047 & 0.87 &  \\
Gemini Flash & 0.029 & 0.017$^{\dagger}$ & 0.015$^{\dagger}$ & 0.046 & 0.037 & 0.62 &  \\
GPT-6 Astra & 0.016 & 0.012$^{\dagger}$ & 0.091 & 0.034 & 0.035 & 0.95 &  \\
\bottomrule
\end{tabular}\end{table}

\begin{table}[t]
\centering
\caption{Pooled procyclicality by prompt framing, notches per severity step with firm-clustered bootstrap 95\% CIs. Pooled across the five non-frontier models, whose four framings form the original confirmatory grid; the two frontier models' framing gradients are reported in the replication archive because their arms were added in revision.}
\label{tab:gradient}
\begin{tabular}{lrrr}
\toprule
Framing & Pooled $\beta$ & 95\% CI & Excl.\ 0 \\
\midrule
Point-in-time & 0.516 & [0.49, 0.54] & Yes \\
Terse & 0.429 & [0.41, 0.45] & Yes \\
"Be stable" & 0.329 & [0.30, 0.35] & Yes \\
Through-the-cycle & 0.160 & [0.14, 0.18] & Yes \\
\bottomrule
\end{tabular}
\end{table}

\begin{figure}[t]
 \centering
 \includegraphics[width=\linewidth]{Figure_1.pdf}
 \caption{(a) Per-model procyclicality $\beta$ under the terse framing, all seven models, with 95\%
 firm-clustered bootstrap CIs. The dashed line marks the through-the-cycle null ($\beta=0$).
 (b) Pooled $\beta$ across the four prompt framings (point-in-time $\to$ terse $\to$ ``be stable''
 $\to$ through-the-cycle) with 95\% CIs, pooled across the five original non-frontier models as in Table~\ref{tab:gradient}. The two frontier arms enter panel~(a) and the per-model tables.}
 \label{fig:main}
\end{figure}

\begin{figure}[t]
 \centering
 \includegraphics[width=0.85\linewidth]{Figure_2.pdf}
 \caption{Asymmetry: per-model downside half-slope (neutral$\to$severe recession) versus upside
 half-slope (boom$\to$neutral), terse framing, with 95\% firm-clustered bootstrap confidence
 intervals. Every model downgrades faster into recessions than it upgrades into booms, and every
 downside-to-upside ratio interval excludes~$1$. The ratio intervals are wide because the upside
 half-slope is small, so the ordering of models by asymmetry is not resolvable. \texttt{(t=1)} marks
 the one arm that could not run at temperature zero.}
 \label{fig:asym}
\end{figure}

\section*{Declaration of generative AI and AI-assisted technologies in the manuscript preparation process}
During the preparation of this work the authors used Claude (Anthropic) to assist in implementing the
data-capture and analysis pipeline and in drafting and editing the manuscript text. The large language
models evaluated in this study, four OpenAI models and three Google Gemini models, constitute the
object of study of this paper, and their use is fully described in Section~\ref{sec:design}. They are
not tools of manuscript preparation. All AI-assisted content was carefully reviewed, edited, and
verified by the authors, who take full responsibility for the content of the published article.

\section*{Declaration of competing interest}
The authors declare that they have no known competing financial interests or personal relationships
that could have appeared to influence the work reported in this paper.

\section*{Funding}
This research did not receive any specific grant from funding agencies in the public, commercial, or
not-for-profit sectors.

\section*{Data availability}
All data and code are openly available in a citable Zenodo archive,
\url{https://doi.org/10.5281/zenodo.21864042} (mirrored on GitHub). Reproduction is offline,
deterministic and free \citep{artifact2026}.

\appendix
\section{Robustness}\label{app:robust}
The original confirmatory grid contains $32{,}000$ assessments. Including the revision experiments,
the study contains $80{,}590$ frozen assessments, every one of which parsed. The rating scale has $21$ points because the default grades D, SD and RD share the bottom notch with C. D accounts for $90$ assessments and SD and RD for none. A further $101$ call attempts, $0.13$\% of the $80{,}691$ made in frozen captures, returned a transport or schema error at capture and are excluded from the frozen set. Table~\ref{tab:firmspec}
reports the synthetic-firm generation specification.

\emph{Placebo and alternative estimators.} The pooled macro slope is $0.4295$ against $0.0425$ for the
matched credit-irrelevant placebo, yielding a placebo-adjusted effect of $0.387$ notches per severity
step (95\% CI $[0.363,0.412]$). The interval is obtained by bootstrapping the slope difference on the
same resampled firms. A joint specification containing both axes returns the same coefficient and the
same interval. Served-model fixed effects shift OpenAI coefficients by less than $0.003$.
Table~\ref{tab:robust} reports the expected-notch, replication-noise and ordered-probit estimates.

\emph{Real fundamentals and issuer recall.} Using FY2024 financial data from $80$ SEC filers, with
names, tickers, CIKs and dates removed before capture, the positive slope replicates in all four
models tested: GPT-4o-mini $0.814$ $[0.729,0.896]$, GPT-4o $0.407$ $[0.362,0.451]$, GPT-6 Astra
$0.151$ $[0.126,0.179]$ and Gemini Flash $0.094$ $[0.069,0.121]$. All intervals exclude zero. Of the
$45$ line-by-tier comparisons between the Table~\ref{tab:firmspec} ranges and the real-firm
interquartile ranges, $41$ overlap. The four exceptions are top-tier revenue and net debt to EBITDA,
and distressed-firm current ratio and working capital to assets. Because synthetic lines are drawn
independently, this supports marginal credit-ratio realism, not firm scale or the joint covariance
structure of real filings. GPT-6 Astra identifies $7/40$ de-identified real firms and GPT-4o $0/40$, versus $0/40$ synthetic controls. The slope difference
between identified and tested-but-unidentified firms is $-0.003$ $[-0.074,0.063]$, with the
identified subsample too small to establish equivalence.

\emph{Replication and model drift.} The primary condition has five replicates per cell for
\texttt{gpt-4o-mini}, \texttt{gpt-4.1-mini}, \texttt{gpt-4o} and \texttt{gemini-flash-latest}, and
two for the remaining models. Three OpenAI arms use dated snapshots
(\texttt{gpt-4o-mini-2024-07-18}, \texttt{gpt-4.1-mini-2025-04-14}, \texttt{gpt-4o-2024-11-20}) and
log \texttt{system\_fingerprint}, making the served backend identifiable per row but not reproducible
later. GPT-6 Astra and the three Gemini \texttt{-latest} endpoints use unresolved aliases.
Cross-epoch coefficient drift remains within sampling uncertainty.

% The last appendix paragraph is long and was breaking mid-sentence at the page foot, with the
% two tables landing in the gap. Starting it on a fresh page gives it room to finish, and the
% \clearpage after it flushes both floats so neither can migrate back into the text.
\clearpage
\emph{Formal contrasts and sensitivity.} Paired firm-clustered bootstraps show that $18$ of $21$
model pairs differ significantly ($16$ after a Holm correction). The ``be stable'' versus through-the-cycle contrast is significant
in all seven models, ranging from $0.047$ $[0.017,0.080]$ to $0.254$ $[0.206,0.304]$. PD-retention
intervals lie below the $0.75$ threshold for every model, ranging from $0.587$ $[0.529,0.642]$ to
$-0.106$ $[-0.166,-0.059]$. Under the ``be stable'' instruction PD retention instead ranges from $0.42$ to $0.92$.

The baseline capital mapping uses a stylised $21$-notch PD ladder, anchored at approximately AAA
$0.01$\%, BBB $0.25$\%, BB $1.1$\%, B $5$\% and CCC $20$\%, with LGD $=0.45$ and maturity $=2.5$
years. The ladder reproduces the ordering and broad magnitudes of the one-year corporate default
rates reported by \citet{spratings2025} rather than copying that study's values. The Basel
$0.05$\% PD input floor of \citet[CRE32.4]{bcbs_cre32} is then applied, which binds on the four
notches at AA$-$ and above. Reported capital figures are on the floored ladder. A $259$-cell
sensitivity grid varies four PD mappings (baseline, PD$\times0.5$, PD$\times2$ and a one-notch downward shift), three LGDs ($0.30$, $0.45$, $0.60$), three maturities
($1.0$, $2.5$, $5.0$ years) and seven models, plus one verbalised-PD mapping per model. The capital
swing remains positive in all $259$ cells. Its magnitude varies by $2.7$ to $4.7$ times across the
rating-based specifications and by $3.1$ to $7.6$ times including verbalised PD. The
$\pm0.05$-notch TTC rating band and $0.75$ PD-retention threshold are reporting conventions rather
than supervisory standards.

\clearpage
\begin{table}[tb]\centering\small
\caption{Synthetic-firm generation specification. Each firm draws every line independently and uniformly from the interval shown for its credit-quality tier; firms are stratified across the five tiers and eight sectors so the fundamentals-implied ratings span the scale. Ranges were chosen so a tier's implied rating lands in its band, not fitted to a reference population. Of the 45 line-by-tier cells, 41 have a draw interval intersecting the interquartile range of the real firms in that tier (Appendix~\ref{app:robust}). Values are frozen and SHA-256'd before capture.}
\label{tab:firmspec}
\begin{tabular}{lrrrrr}
\toprule
Line & High IG & Low IG & Crossover & High yield & Distressed \\
\midrule
Revenue (USD m) & 40000 to 180000 & 8000 to 60000 & 2000 to 20000 & 500 to 8000 & 100 to 3000 \\
Revenue growth (\%) & 2 to 9 & 1 to 7 & $-$2 to 6 & $-$5 to 6 & $-$18 to 2 \\
EBITDA margin (\%) & 22 to 38 & 15 to 27 & 11 to 20 & 7 to 16 & 0 to 9 \\
Net debt / EBITDA ($\times$) & 0.2 to 1.3 & 1.3 to 2.8 & 2.8 to 4.2 & 4.2 to 6.5 & 6.5 to 12.0 \\
EBITDA / interest ($\times$) & 14.0 to 40.0 & 6.0 to 13.0 & 3.0 to 6.0 & 1.6 to 3.2 & 0.4 to 1.6 \\
FCF margin (\%) & 9 to 18 & 4 to 10 & 0 to 6 & $-$4 to 3 & $-$14 to $-$1 \\
Current ratio & 1.4 to 2.4 & 1.1 to 1.8 & 0.9 to 1.4 & 0.7 to 1.2 & 0.4 to 0.9 \\
Retained earnings / assets & 0.35 to 0.60 & 0.18 to 0.40 & 0.05 to 0.22 & $-$0.05 to 0.10 & $-$0.40 to $-$0.02 \\
Working capital / assets & 0.15 to 0.32 & 0.06 to 0.20 & $-$0.02 to 0.12 & $-$0.08 to 0.06 & $-$0.25 to $-$0.02 \\
\bottomrule
\end{tabular}\end{table}

\begin{table}[t]
\centering
\caption{Robustness of the terse-framing coefficient. Net $=$ placebo-adjusted coefficient. Exp-n. $=$ expected-notch slope from rating-token log-probabilities. Noise SD $=$ mean within-cell notch standard deviation across replicate seeds. Probit $=$ Amato--Furfine ordered-probit estimate, converted from the latent scale to notches per step. Expected-notch estimates are unavailable for Gemini and for GPT-6 Astra because neither returns rating-token log-probabilities. Gemini Flash-Lite's $0.000$ Noise SD is exact.}
\label{tab:robust}
\begin{tabular}{lrrrrrr}
\toprule
Model & $\beta$ & Net & Exp-n. terse & Exp-n. TTC & Noise SD & Probit \\
\midrule
Gemini Pro & 0.182 & 0.174 &  &  & 0.080 & 0.138 \\
GPT-6 Astra & 0.198 & 0.186 &  &  & 0.067 & 0.177 \\
Gemini Flash & 0.230 & 0.190 &  &  & 0.020 & 0.272 \\
Gemini Flash-Lite & 0.285 & 0.268 &  &  & 0.000 & 0.235 \\
GPT-4o & 0.387 & 0.367 & 0.378 & 0.182 & 0.025 & 0.403 \\
GPT-4.1-mini & 0.471 & 0.464 & 0.420 & 0.081 & 0.086 & 0.307 \\
GPT-4o-mini & 0.775 & 0.646 & 0.765 & 0.418 & 0.072 & 0.726 \\
\bottomrule
\end{tabular}
\end{table}

\section{Prompts}\label{app:prompts}
The prompt framings and the macro and placebo blocks below are extracted programmatically from the
capture code (\texttt{agent.py}) and reproduce the executed instrument verbatim.
% AUTO-GENERATED from capture/agent.py by analysis/make_prompts.py — do not edit.

\subsection{Prompt framings}\label{app:framings}

The four instruction framings (system prompt), verbatim as executed:

\paragraph{\texttt{v\_terse}}

\begin{lstlisting}[breaklines=true,breakatwhitespace=false,
                   basicstyle=\ttfamily\scriptsize,columns=fullflexible,
                   frame=none,xleftmargin=0pt,literate={\%}{{\%}}1]
You are a credit rating analyst. Assign a long-term issuer credit rating to the company described, based on its fundamentals and the stated operating environment. Respond with ONLY a JSON object and nothing else, in exactly this form: {"rating": "<one S&P long-term issuer rating from AAA down to D, using +/- modifiers, e.g. AA-, BBB+, B>", "pd_1yr_bps": <your estimated 1-year probability of default in BASIS POINTS as an integer, where 100 bps = 1%>, "rationale": "<=30 words"}.
\end{lstlisting}

\paragraph{\texttt{v\_ttc}}

\begin{lstlisting}[breaklines=true,breakatwhitespace=false,
                   basicstyle=\ttfamily\scriptsize,columns=fullflexible,
                   frame=none,xleftmargin=0pt,literate={\%}{{\%}}1]
You are a credit rating analyst applying a THROUGH-THE-CYCLE methodology. Assign a long-term issuer credit rating that reflects the company's creditworthiness over a full economic cycle, deliberately LOOKING THROUGH transitory macroeconomic conditions rather than the current point in the cycle. Respond with ONLY a JSON object and nothing else, in exactly this form: {"rating": "<one S&P long-term issuer rating from AAA down to D, using +/- modifiers, e.g. AA-, BBB+, B>", "pd_1yr_bps": <your estimated 1-year probability of default in BASIS POINTS as an integer, where 100 bps = 1%>, "rationale": "<=30 words"}.
\end{lstlisting}

\paragraph{\texttt{v\_pit}}

\begin{lstlisting}[breaklines=true,breakatwhitespace=false,
                   basicstyle=\ttfamily\scriptsize,columns=fullflexible,
                   frame=none,xleftmargin=0pt,literate={\%}{{\%}}1]
You are a credit rating analyst applying a POINT-IN-TIME methodology. Assign a long-term issuer credit rating that reflects the company's creditworthiness given CURRENT economic and financial conditions as described. Respond with ONLY a JSON object and nothing else, in exactly this form: {"rating": "<one S&P long-term issuer rating from AAA down to D, using +/- modifiers, e.g. AA-, BBB+, B>", "pd_1yr_bps": <your estimated 1-year probability of default in BASIS POINTS as an integer, where 100 bps = 1%>, "rationale": "<=30 words"}.
\end{lstlisting}

\paragraph{\texttt{v\_stable}}

\begin{lstlisting}[breaklines=true,breakatwhitespace=false,
                   basicstyle=\ttfamily\scriptsize,columns=fullflexible,
                   frame=none,xleftmargin=0pt,literate={\%}{{\%}}1]
You are a credit rating analyst. Assign a long-term issuer credit rating that is STABLE and does not change much in response to short-term conditions. Respond with ONLY a JSON object and nothing else, in exactly this form: {"rating": "<one S&P long-term issuer rating from AAA down to D, using +/- modifiers, e.g. AA-, BBB+, B>", "pd_1yr_bps": <your estimated 1-year probability of default in BASIS POINTS as an integer, where 100 bps = 1%>, "rationale": "<=30 words"}.
\end{lstlisting}

\subsection{Macro-environment blocks}\label{app:macro}

Prepended to the byte-identical fundamentals; only these five blocks differ across the severity axis:

\paragraph{boom}

\begin{lstlisting}[breaklines=true,breakatwhitespace=false,
                   basicstyle=\ttfamily\scriptsize,columns=fullflexible,
                   frame=none,xleftmargin=0pt,literate={\%}{{\%}}1]
MACRO ENVIRONMENT (current conditions, this fiscal period):
  - Real GDP growth (annualized): +4.5%
  - Unemployment rate: 3.4%
  - Investment-grade corporate credit spread: 95 bps
  - Trailing 12-month corporate default rate: 0.6%
  - Bank lending standards: easing broadly
\end{lstlisting}

\paragraph{expansion}

\begin{lstlisting}[breaklines=true,breakatwhitespace=false,
                   basicstyle=\ttfamily\scriptsize,columns=fullflexible,
                   frame=none,xleftmargin=0pt,literate={\%}{{\%}}1]
MACRO ENVIRONMENT (current conditions, this fiscal period):
  - Real GDP growth (annualized): +2.6%
  - Unemployment rate: 4.2%
  - Investment-grade corporate credit spread: 135 bps
  - Trailing 12-month corporate default rate: 1.4%
  - Bank lending standards: modestly easing
\end{lstlisting}

\paragraph{neutral}

\begin{lstlisting}[breaklines=true,breakatwhitespace=false,
                   basicstyle=\ttfamily\scriptsize,columns=fullflexible,
                   frame=none,xleftmargin=0pt,literate={\%}{{\%}}1]
MACRO ENVIRONMENT (current conditions, this fiscal period):
  - Real GDP growth (annualized): +1.8%
  - Unemployment rate: 5.0%
  - Investment-grade corporate credit spread: 175 bps
  - Trailing 12-month corporate default rate: 2.3%
  - Bank lending standards: unchanged
\end{lstlisting}

\paragraph{slowdown}

\begin{lstlisting}[breaklines=true,breakatwhitespace=false,
                   basicstyle=\ttfamily\scriptsize,columns=fullflexible,
                   frame=none,xleftmargin=0pt,literate={\%}{{\%}}1]
MACRO ENVIRONMENT (current conditions, this fiscal period):
  - Real GDP growth (annualized): -0.8%
  - Unemployment rate: 6.4%
  - Investment-grade corporate credit spread: 320 bps
  - Trailing 12-month corporate default rate: 4.1%
  - Bank lending standards: tightening
\end{lstlisting}

\paragraph{severe\_recession}

\begin{lstlisting}[breaklines=true,breakatwhitespace=false,
                   basicstyle=\ttfamily\scriptsize,columns=fullflexible,
                   frame=none,xleftmargin=0pt,literate={\%}{{\%}}1]
MACRO ENVIRONMENT (current conditions, this fiscal period):
  - Real GDP growth (annualized): -3.6%
  - Unemployment rate: 8.9%
  - Investment-grade corporate credit spread: 620 bps
  - Trailing 12-month corporate default rate: 7.5%
  - Bank lending standards: sharply tightening
\end{lstlisting}

\subsection{Placebo blocks (credit-irrelevant control)}\label{app:placebo}

Matched in structure and number density, on a credit-irrelevant axis:

\paragraph{placebo\_calm}

\begin{lstlisting}[breaklines=true,breakatwhitespace=false,
                   basicstyle=\ttfamily\scriptsize,columns=fullflexible,
                   frame=none,xleftmargin=0pt,literate={\%}{{\%}}1]
REGIONAL CONDITIONS (current, at the company's headquarters region):
  - Average temperature vs seasonal norm: +0.6 degC
  - Daylight hours vs annual mean: +4.5%
  - Monthly rainfall vs normal: 95 mm
  - Regional air-quality index: 38
  - Commuter traffic congestion: light
\end{lstlisting}

\paragraph{placebo\_mild}

\begin{lstlisting}[breaklines=true,breakatwhitespace=false,
                   basicstyle=\ttfamily\scriptsize,columns=fullflexible,
                   frame=none,xleftmargin=0pt,literate={\%}{{\%}}1]
REGIONAL CONDITIONS (current, at the company's headquarters region):
  - Average temperature vs seasonal norm: +0.2 degC
  - Daylight hours vs annual mean: +2.6%
  - Monthly rainfall vs normal: 135 mm
  - Regional air-quality index: 52
  - Commuter traffic congestion: moderate
\end{lstlisting}

\paragraph{placebo\_normal}

\begin{lstlisting}[breaklines=true,breakatwhitespace=false,
                   basicstyle=\ttfamily\scriptsize,columns=fullflexible,
                   frame=none,xleftmargin=0pt,literate={\%}{{\%}}1]
REGIONAL CONDITIONS (current, at the company's headquarters region):
  - Average temperature vs seasonal norm: 0.0 degC
  - Daylight hours vs annual mean: +1.8%
  - Monthly rainfall vs normal: 175 mm
  - Regional air-quality index: 60
  - Commuter traffic congestion: typical
\end{lstlisting}

\paragraph{placebo\_rough}

\begin{lstlisting}[breaklines=true,breakatwhitespace=false,
                   basicstyle=\ttfamily\scriptsize,columns=fullflexible,
                   frame=none,xleftmargin=0pt,literate={\%}{{\%}}1]
REGIONAL CONDITIONS (current, at the company's headquarters region):
  - Average temperature vs seasonal norm: -0.4 degC
  - Daylight hours vs annual mean: -0.8%
  - Monthly rainfall vs normal: 320 mm
  - Regional air-quality index: 88
  - Commuter traffic congestion: heavy
\end{lstlisting}

\paragraph{placebo\_storm}

\begin{lstlisting}[breaklines=true,breakatwhitespace=false,
                   basicstyle=\ttfamily\scriptsize,columns=fullflexible,
                   frame=none,xleftmargin=0pt,literate={\%}{{\%}}1]
REGIONAL CONDITIONS (current, at the company's headquarters region):
  - Average temperature vs seasonal norm: -1.1 degC
  - Daylight hours vs annual mean: -3.6%
  - Monthly rainfall vs normal: 620 mm
  - Regional air-quality index: 140
  - Commuter traffic congestion: severe
\end{lstlisting}

\printcredits

\bibliographystyle{cas-model2-names}
\bibliography{refs}

\end{document}
